\documentclass[11pt]{article}

\usepackage{amsmath,amssymb,amsthm,mathtools}
\usepackage[margin=1.02in]{geometry}
\usepackage{hyperref}

\newtheorem{theorem}{Theorem}
\newtheorem{lemma}{Lemma}
\newtheorem{corollary}{Corollary}
\newtheorem{proposition}{Proposition}
\newtheorem{definition}{Definition}
\newtheorem{remark}{Remark}

\newcommand{\E}{\mathbb{E}}
\newcommand{\cC}{\mathcal{C}}
\newcommand{\cP}{\mathcal{P}}
\newcommand{\cN}{\mathcal{N}}
\newcommand{\Tail}{\operatorname{Tail}}
\newcommand{\supp}{\operatorname{supp}}

\title{The Chromatic-Polynomial Graphon Bound for Pure Chordal Graphs}
\author{}
\date{}

\begin{document}

\maketitle

\begin{abstract}
We give a graphon-direct proof of the chromatic-polynomial lower bound for
chordal graphs whose maximal cliques have one common size.  Every analytic
inequality is stated in denominator-free form until positivity has been
proved.  The proof uses only finite-product integration, an elementary Gibbs
inequality, the regularization
\(W_\varepsilon=\varepsilon+(1-\varepsilon)W\), a direct graphon proof of the
Moon--Moser recurrence, and finite clique-tree combinatorics.  It uses no
finite-graph approximation, graphon regularity, cut metric, injective-copy
count, minimum-degree assumption, or pointwise degree lower bound.

The organization and dependency boundaries mirror the machine-checked Lean
development in \path{lean/Taeyoung/Methods/PureChordal}.  In particular, the
primary analytic result is a polynomial certificate inequality; the
chromatic-polynomial identity is proved separately.
\end{abstract}

\section{Graphons and the theorem}

Let \((\Omega,\mu)\) be a probability space.  A graphon is a symmetric
measurable function
\[
    W\colon\Omega^2\longrightarrow[0,1].
\]
For a finite simple graph \(F\), define
\[
    t(F,W)
    :=
    \int_{\Omega^{V(F)}}
       \prod_{uv\in E(F)}W(x_u,x_v)
       \prod_{u\in V(F)}d\mu(x_u).
\]
All integrands are measurable, nonnegative, and at most \(1\), so every
displayed integral exists and belongs to \([0,1]\).  Write
\[
    p:=t(K_2,W),
    \qquad
    t_s:=t(K_s,W),
    \qquad
    t_0=t_1:=1.
\]

For an integer \(s\ge0\), put
\[
    A_s(p)
    :=
    \prod_{a=0}^{s-1}\bigl(1-a(1-p)\bigr),
    \qquad A_0(p)=A_1(p)=1.
\]
For \(0\le s\le r\), also put
\[
    \Tail_{s,r}(p)
    :=
    \prod_{a=s}^{r-1}\bigl(1-a(1-p)\bigr).
\]
The empty product is \(1\), and the exact factorization
\begin{equation}\label{eq:A-tail}
    A_s(p)\Tail_{s,r}(p)=A_r(p)
\end{equation}
holds as a polynomial identity, without any division.

\begin{theorem}[Pure-chordal graphon inequality]\label{thm:main}
Let \(H\) be a finite chordal graph all of whose maximal cliques have the
same cardinality \(r\ge3\).  Connectivity is not required.  Then
\[
    \chi(H)=r.
\]
For every graphon \(W\) of edge density \(p\) satisfying
\[
    p\ge1-\frac1{r-1},
\]
one has, for \(p<1\),
\[
    \boxed{
    t(H,W)
    \ge
    (1-p)^{v(H)}
    \chi_H\!\left(\frac1{1-p}\right).
    }
\]
At \(p=1\), the right-hand side is interpreted by its polynomial
continuation in \(p\); that continuation equals \(1\), and the same inequality
holds with equality.
\end{theorem}

The remainder of the paper proves every component of this statement.

\section{The rooted pure clique-tree certificate}

The analytic proof uses explicit finite certificate data.  This avoids hiding
the running-intersection property inside informal clique-tree language.

\begin{definition}[Rooted pure clique-tree presentation]\label{def:certificate}
Let \(m\ge1\).  A rooted \(r\)-pure clique-tree presentation of \(H\) consists
of distinct bags
\[
    C_0,C_1,\dots,C_{m-1}\subseteq V(H)
\]
and, for each \(i>0\), a parent index \(\pi(i)<i\), such that:
\begin{enumerate}
    \item every \(C_i\) has cardinality \(r\) and induces a clique in \(H\);
    \item the bags cover every vertex and every edge of \(H\);
    \item for every \(i>0\),
    \begin{equation}\label{eq:old-separator}
        C_i\cap\bigcup_{j<i}C_j
        =
        C_i\cap C_{\pi(i)}.
    \end{equation}
\end{enumerate}
Set
\[
    S_0:=\varnothing,
    \qquad
    S_i:=C_i\cap C_{\pi(i)}\quad(i>0),
    \qquad
    s_i:=|S_i|,
\]
and let
\[
    N_i:=C_i\setminus S_i
\]
be the vertices introduced at step \(i\).
\end{definition}

For a chordal graph, the standard maximal-clique-tree theorem supplies a tree
on the maximal cliques with the running-intersection property.  Rooting that
tree and listing parents before children gives
Definition~\ref{def:certificate}.  For a disconnected chordal graph, take a
clique tree in each component and join those trees by artificial tree edges;
their separators are empty.  Thus Theorem~\ref{thm:main} reduces to the
certificate theorem proved below.  Conversely, the Lean development uses the
existence of this maximal-clique presentation as its formal definition of
chordality, so no unrecorded graph-theoretic implication occurs in the formal
theorem.

\begin{lemma}[Strict separator size]\label{lem:separator-size}
For every bag index \(i\),
\[
    0\le s_i<r.
\]
\end{lemma}

\begin{proof}
For the root, \(s_0=0<r\).  If \(i>0\), then \(S_i\subseteq C_i\), so
\(s_i\le r\).  Equality would give \(S_i=C_i\), hence
\(C_i\subseteq C_{\pi(i)}\).  The parent also has \(r\) vertices, so this
would force \(C_i=C_{\pi(i)}\), contradicting distinctness of the bags.
\end{proof}

\begin{lemma}[Introduced-vertex partition]\label{lem:new-partition}
The sets \(N_0,\dots,N_{m-1}\) are pairwise disjoint and their union is
\(V(H)\).  Consequently,
\begin{equation}\label{eq:vertex-count}
    v(H)
    =
    \sum_{i=0}^{m-1}(r-s_i)
    =
    mr-\sum_{i=0}^{m-1}s_i.
\end{equation}
\end{lemma}

\begin{proof}
Let \(v\in V(H)\), and choose the least index \(i\) for which \(v\in C_i\).
Then \(v\notin\bigcup_{j<i}C_j\), so
\eqref{eq:old-separator} implies \(v\notin S_i\).  Hence \(v\in N_i\), proving
that the \(N_i\) cover \(V(H)\).

If \(i<j\) and \(v\in N_i\cap N_j\), then \(v\in C_i\) shows that \(v\) is
already present before step \(j\).  Since also \(v\in C_j\),
\eqref{eq:old-separator} gives \(v\in S_j\), contrary to \(v\in N_j\).
Thus the sets are pairwise disjoint.  Finally,
\(|N_i|=|C_i|-|S_i|=r-s_i\), and summing the cardinalities proves
\eqref{eq:vertex-count}.
\end{proof}

For a finite set \(A\subseteq V(H)\), let \(\binom A2\) be the set of
unordered pairs of distinct elements of \(A\).  For an assignment
\(x\in\Omega^{V(H)}\), define
\[
    \kappa_A(x)
    :=
    \prod_{\{u,v\}\in\binom A2}W(x_u,x_v),
    \qquad
    \kappa_H(x)
    :=
    \prod_{uv\in E(H)}W(x_u,x_v).
\]

\begin{lemma}[Edge-weight multiplicity]\label{lem:edge-multiplicity}
For every graphon \(W\) and every assignment \(x\),
\begin{equation}\label{eq:edge-weight}
    \prod_{i=0}^{m-1}\kappa_{C_i}(x)
    =
    \kappa_H(x)
    \prod_{i=0}^{m-1}\kappa_{S_i}(x).
\end{equation}
\end{lemma}

\begin{proof}
First observe that
\begin{equation}\label{eq:pairs-overlap}
    \binom{C_i}{2}
    \cap
    \bigcup_{j<i}\binom{C_j}{2}
    =
    \binom{S_i}{2}.
\end{equation}
For the nontrivial inclusion, if both endpoints of a pair lie in \(C_i\) and
also in an earlier bag, then each endpoint lies in
\(C_i\cap\bigcup_{j<i}C_j=S_i\).  Conversely, when \(i>0\), every pair in
\(S_i\) also lies in the earlier parent bag; for \(i=0\), both sides are
empty.

Now insert the bags one at a time.  For any commuting family of weights,
the product over a set \(A\), times the product over a set \(B\), is the
product over \(A\cup B\), times the product over \(A\cap B\).  Applying this
identity at step \(i\) and using \eqref{eq:pairs-overlap} shows inductively
that the product of all bag-pair weights equals the product over the union of
all bag pairs, times the product of the separator-pair weights.  The union of
the bag pairs is exactly \(E(H)\): bag cliquehood gives one inclusion, and
edge coverage gives the other.  This is precisely \eqref{eq:edge-weight}.
\end{proof}

\section{Chromatic polynomial and the exact target}

Write
\[
    (x)_d:=x(x-1)\cdots(x-d+1),
    \qquad (x)_0:=1.
\]

\begin{lemma}[Chromatic factorization]\label{lem:chromatic}
For a graph with a rooted \(r\)-pure clique-tree presentation,
\begin{equation}\label{eq:chromatic-product}
    \chi_H(x)
    =
    \prod_{i=0}^{m-1}(x-s_i)_{r-s_i}.
\end{equation}
Equivalently,
\begin{equation}\label{eq:chromatic-quotient}
    \chi_H(x)
    =
    \frac{(x)_r^m}{\prod_{i=0}^{m-1}(x)_{s_i}},
\end{equation}
where \eqref{eq:chromatic-product}, not division, is the primary polynomial
identity.  Moreover, \(\chi(H)=r\).
\end{lemma}

\begin{proof}
Fix an integer \(k\ge r\) and expose bags in index order.  A completed
assignment is a proper \(k\)-colouring of \(H\) if and only if its restriction
to every bag is injective.  Indeed, properness implies injectivity on each
clique, while injectivity on all bags separates the endpoints of every edge
because every edge lies in a bag.

At step \(i\), the already exposed vertices of \(C_i\) are exactly \(S_i\).
Their \(s_i\) colours are distinct because, for \(i>0\), they lie in the
parent bag.  The \(r-s_i\) vertices of \(N_i\) must receive distinct colours
from the remaining \(k-s_i\) colours, which can be done in exactly
\[
    (k-s_i)_{r-s_i}
\]
ways.  This choice makes the current bag injective.  It cannot create a
future obstruction: the old vertices of every future bag form its separator
inside an already injectively coloured parent bag.  Induction therefore gives
\[
    \chi_H(k)
    =
    \prod_{i=0}^{m-1}(k-s_i)_{r-s_i}
    \qquad(k\ge r).
\]
Both sides are polynomials in \(k\), and they agree at infinitely many
integers, proving \eqref{eq:chromatic-product}.  The identity
\((x-s)_{r-s}(x)_s=(x)_r\) gives the cross-multiplied form of
\eqref{eq:chromatic-quotient}.

Every bag contains a copy of \(K_r\), so \(\chi(H)\ge r\).  At \(k=r\), every
factor \((r-s_i)_{r-s_i}\) is positive by
Lemma~\ref{lem:separator-size}; hence \(H\) has an \(r\)-colouring.  Thus
\(\chi(H)=r\).
\end{proof}

\begin{lemma}[Polynomial certificate equals the target]\label{lem:target}
Define
\begin{equation}\label{eq:certificate}
    B_H(p)
    :=
    \prod_{i=0}^{m-1}\Tail_{s_i,r}(p).
\end{equation}
For \(p\ne1\),
\begin{equation}\label{eq:target-identity}
    B_H(p)
    =
    (1-p)^{v(H)}
    \chi_H\!\left(\frac1{1-p}\right).
\end{equation}
Thus \(B_H\) is the polynomial continuation of the target to \(p=1\), and
\[
    B_H(1)=1.
\]
Also,
\begin{equation}\label{eq:certificate-cross}
    \left(\prod_{i=0}^{m-1}A_{s_i}(p)\right)B_H(p)
    =
    A_r(p)^m.
\end{equation}
\end{lemma}

\begin{proof}
Put \(q=1-p\ne0\).  For every \(s\le r\),
\[
\begin{aligned}
    q^{r-s}\left(\frac1q-s\right)_{r-s}
    &=
    \prod_{a=s}^{r-1}q\left(\frac1q-a\right)\\
    &=
    \prod_{a=s}^{r-1}(1-aq)
    =\Tail_{s,r}(p).
\end{aligned}
\]
Multiply this identity over all bags.  The exponent of \(q\) is
\(\sum_i(r-s_i)=v(H)\) by \eqref{eq:vertex-count}, and
Lemma~\ref{lem:chromatic} supplies the remaining product.  This proves
\eqref{eq:target-identity}.  At \(p=1\), every factor in
\eqref{eq:certificate} equals \(1\).  Finally, multiplying
\eqref{eq:A-tail} over all \(i\) gives
\eqref{eq:certificate-cross}.
\end{proof}

\section{Clique-tree gluing directly in graphon space}

The main analytic statement in this section is the cross-multiplied inequality
\begin{equation}\label{eq:gluing-goal}
    t_r^m
    \le
    t(H,W)\prod_{i=0}^{m-1}t_{s_i}.
\end{equation}
Because \(s_0=0\) and \(t_0=1\), the root factor is harmless.  The
cross-multiplied form remains meaningful when a clique density is zero.

\subsection{A local Gibbs lemma}

\begin{lemma}[Gibbs inequality]\label{lem:gibbs}
Let \(f,g\) be strictly positive measurable probability densities on a
probability space.  Suppose \(f,g\), their reciprocals, and the logarithmic
integrand below are bounded.  Then
\[
    \int f\log\frac{g}{f}\,d\mu\le0.
\]
\end{lemma}

\begin{proof}
For \(u>0\), \(\log u\le u-1\).  Substitute \(u=g/f\), multiply by
the positive function \(f\), and integrate:
\[
    \int f\log\frac{g}{f}
    \le
    \int(g-f)
    =1-1=0.
\]
The boundedness assumptions justify integrability; in the application they
follow from a uniform positive lower bound on \(W\).
\end{proof}

\subsection{Uniformly positive graphons}

Assume temporarily that there is a number \(\delta>0\) such that
\begin{equation}\label{eq:positive-W}
    \delta\le W(x,y)\le1
    \qquad\text{for every }x,y\in\Omega.
\end{equation}
Every finite graph weight and every clique density is then strictly positive.

For each bag \(C_i\), define its unnormalized marginal over the new
coordinates by
\[
    R_i(x_{S_i})
    :=
    \int_{\Omega^{N_i}}
       \kappa_{C_i}(x_{S_i},y_{N_i})
       \prod_{v\in N_i}d\mu(y_v).
\]
For the root this is the constant \(R_0=t_r\).  Define
\[
    q_i(x_{S_i}):=\frac{R_i(x_{S_i})}{t_r},
    \qquad
    K_i(y_{N_i}\mid x_{S_i})
    :=
    \frac{\kappa_{C_i}(x_{S_i},y_{N_i})}{R_i(x_{S_i})}.
\]
Then \(q_i\) is the separator marginal of the normalized \(K_r\)-weight,
and, for every fixed \(x_{S_i}\),
\begin{equation}\label{eq:conditional-one}
    \int_{\Omega^{N_i}}K_i(y_{N_i}\mid x_{S_i})\,d\mu^{N_i}(y)=1.
\end{equation}
All quantities are measurable by finite-product Fubini/Tonelli.  Moreover,
\eqref{eq:positive-W} bounds every numerator and denominator between explicit
positive powers of \(\delta\) and \(1\), so all quotients and logarithms used
below are bounded.

\begin{lemma}[Equal-size clique marginal]\label{lem:equal-marginal}
Let \(C,D\) be two \(r\)-element sets and let \(S\subseteq C\cap D\), with
\(C\cap D=S\).  The marginal on the coordinates \(S\) of the normalized
clique weight \(\kappa_C/t_r\) is equal to the corresponding marginal of
\(\kappa_D/t_r\).
\end{lemma}

\begin{proof}
Choose a bijection from \(C\setminus S\) to \(D\setminus S\), and extend it by
the identity on \(S\).  This is possible because both complements have
cardinality \(r-|S|\).  Under the resulting coordinate relabeling, the
complete set of unordered pairs in \(C\) is carried bijectively to the
complete set of unordered pairs in \(D\).  Symmetry of \(W\) identifies the
two integrands, while invariance of a finite product measure under coordinate
permutation identifies the integrals over the complementary coordinates.
The normalizing mass is \(t_r\) for both cliques.
\end{proof}

Define the completed junction density on \(\Omega^{V(H)}\) by
\begin{equation}\label{eq:junction}
    P(x)
    :=
    \prod_{i=0}^{m-1}K_i(x_{N_i}\mid x_{S_i}).
\end{equation}

\begin{lemma}[Junction invariant]\label{lem:junction}
The function \(P\) is a strictly positive probability density.  For every
bag \(C_i\), its \(C_i\)-marginal is
\[
    \frac{\kappa_{C_i}}{t_r},
\]
and for every separator \(S_i\), its \(S_i\)-marginal is \(q_i\).
\end{lemma}

\begin{proof}
Consider the partial products in \eqref{eq:junction}, adding the factors in
the order \(0,1,\dots,m-1\).  We prove inductively that the partial product
has total mass one and that every already inserted bag has normalized
clique-weight marginal.

At the root, \(S_0=\varnothing\), \(R_0=t_r\), and
\(K_0=\kappa_{C_0}/t_r\), so the claim is immediate.

Suppose now that bags before \(i\) have been inserted.  The variables in
\(N_i\) are fresh by Lemma~\ref{lem:new-partition}.  Consequently,
integrating the new factor over \(N_i\) removes it by
\eqref{eq:conditional-one}; this preserves total mass and every earlier bag
marginal.  If \(i>0\), the separator \(S_i\) lies in the parent bag, whose
marginal is \(\kappa_{C_{\pi(i)}}/t_r\) by induction.  By
Lemma~\ref{lem:equal-marginal}, its marginal on \(S_i\) is exactly
\(q_i=R_i/t_r\).  For the root, the same statement is \(1=1\).
Multiplication by \(K_i=\kappa_{C_i}/R_i\) therefore gives the new bag
marginal
\[
    q_i\frac{\kappa_{C_i}}{R_i}
    =
    \frac{R_i}{t_r}\frac{\kappa_{C_i}}{R_i}
    =
    \frac{\kappa_{C_i}}{t_r}.
\]
This completes the induction.  Taking its separator marginal gives \(q_i\).
Strict positivity follows from \eqref{eq:positive-W}.
\end{proof}

For every \(i\), define the normalized separator clique density
\[
    \nu_i(x_{S_i})
    :=
    \frac{\kappa_{S_i}(x_{S_i})}{t_{s_i}}.
\]
The root case is \(\nu_0=1\).  Lift this separator law to a density on all
assignments by
\[
    G_i(x)
    :=
    P(x)\frac{\nu_i(x_{S_i})}{q_i(x_{S_i})}.
\]
Lemma~\ref{lem:junction} shows that \(\int G_i=1\): integrating outside
\(S_i\) replaces \(P\) by \(q_i\), after which the quotient cancels and
\(\int\nu_i=1\).  Also define the normalized graph-weight density
\[
    Q_H(x):=\frac{\kappa_H(x)}{t(H,W)}.
\]
All of \(P,G_i,Q_H\) are strictly positive bounded probability densities.

\begin{lemma}[Constant Gibbs-ratio identity]\label{lem:constant-ratio}
Under \eqref{eq:positive-W}, for every assignment \(x\),
\begin{equation}\label{eq:constant-ratio}
    \frac{Q_H(x)}{P(x)}
    \prod_{i=0}^{m-1}\frac{G_i(x)}{P(x)}
    =
    \frac{t_r^m}
    {t(H,W)\prod_{i=0}^{m-1}t_{s_i}}.
\end{equation}
\end{lemma}

\begin{proof}
By definition,
\[
    P(x)=\frac{\prod_i\kappa_{C_i}(x)}{\prod_iR_i(x_{S_i})}.
\]
Lemma~\ref{lem:edge-multiplicity} changes the numerator to
\(\kappa_H\prod_i\kappa_{S_i}\).  Furthermore,
\[
    \frac{G_i}{P}
    =
    \frac{\nu_i}{q_i}
    =
    \frac{\kappa_{S_i}t_r}{t_{s_i}R_i}.
\]
Substituting these formulas into the left-hand side of
\eqref{eq:constant-ratio} cancels, pointwise, \(\kappa_H\), every
\(\kappa_{S_i}\), and every \(R_i\).  What remains is exactly the constant
on the right-hand side.
\end{proof}

\begin{proposition}[Gluing for a uniformly positive graphon]
Under \eqref{eq:positive-W}, inequality \eqref{eq:gluing-goal} holds.
\end{proposition}

\begin{proof}
Apply Lemma~\ref{lem:gibbs} once to \(f=P,g=Q_H\), and once for each \(i\)
to \(f=P,g=G_i\).  Adding the resulting nonpositive quantities gives
\[
\begin{aligned}
    0
    &\ge
    \int P\log\frac{Q_H}{P}
    +
    \sum_{i=0}^{m-1}\int P\log\frac{G_i}{P}\\
    &=
    \int P\log\left(
      \frac{Q_H}{P}\prod_i\frac{G_i}{P}
    \right).
\end{aligned}
\]
By Lemma~\ref{lem:constant-ratio}, the logarithm is constant, and
\(\int P=1\).  Therefore
\[
    \log\left(
      \frac{t_r^m}{t(H,W)\prod_i t_{s_i}}
    \right)
    \le0.
\]
The argument of the logarithm is positive, so it is at most \(1\).  This is
exactly \eqref{eq:gluing-goal}.
\end{proof}

\subsection{Removing strict positivity}

\begin{lemma}[Finite-product difference]\label{lem:product-difference}
If \(a_e,b_e\in[0,1]\) for every element of a finite set \(E\), then
\[
    \left|\prod_{e\in E}a_e-\prod_{e\in E}b_e\right|
    \le
    \sum_{e\in E}|a_e-b_e|.
\]
\end{lemma}

\begin{proof}
Order the finite set and replace the factors one at a time.  The resulting
telescoping sum has one factor \(a_e-b_e\) in each term, while all other
factors lie in \([0,1]\).  The triangle inequality gives the result.
\end{proof}

For \(0<\varepsilon\le1\), define
\[
    W_\varepsilon
    :=
    \varepsilon+(1-\varepsilon)W.
\]
Then \(\varepsilon\le W_\varepsilon\le1\) pointwise and
\(|W_\varepsilon-W|\le\varepsilon\).  Applying
Lemma~\ref{lem:product-difference} inside the homomorphism-density integral
gives, for every finite graph \(F\),
\begin{equation}\label{eq:regularization-error}
    |t(F,W_\varepsilon)-t(F,W)|
    \le
    e(F)\varepsilon.
\end{equation}

\begin{theorem}[Cross-multiplied clique-tree gluing]\label{thm:gluing}
For every graphon \(W\), with no positivity assumption,
\[
    t_r^m
    \le
    t(H,W)\prod_{i=0}^{m-1}t_{s_i}.
\]
\end{theorem}

\begin{proof}
Let
\[
    L(W):=t_r^m,
    \qquad
    R(W):=t(H,W)\prod_i t_{s_i}.
\]
All densities in these products lie in \([0,1]\).  Applying
Lemma~\ref{lem:product-difference} to the \(m\) identical factors in \(L\),
and then using \eqref{eq:regularization-error} for \(K_r\), gives
\[
    |L(W_\varepsilon)-L(W)|
    \le
    m\binom r2\varepsilon.
\]
Applying the same product lemma to \(R\) gives
\[
    |R(W_\varepsilon)-R(W)|
    \le
    \left(e(H)+\sum_i\binom{s_i}{2}\right)\varepsilon.
\]
Set
\[
    C:=m\binom r2+e(H)+\sum_i\binom{s_i}{2}.
\]
For every \(\eta>0\), choose
\[
    0<\varepsilon\le
    \min\left\{1,\frac{\eta}{C+1}\right\}.
\]
The uniformly positive case applies to \(W_\varepsilon\), so
\(L(W_\varepsilon)\le R(W_\varepsilon)\).  Hence
\[
\begin{aligned}
    L(W)
    &\le L(W_\varepsilon)+m\binom r2\varepsilon\\
    &\le R(W_\varepsilon)+m\binom r2\varepsilon\\
    &\le R(W)+C\varepsilon\\
    &\le R(W)+\eta.
\end{aligned}
\]
Since this holds for every \(\eta>0\), \(L(W)\le R(W)\).  This is an
elementary epsilon argument, not an invocation of graphon approximation or
cut-distance continuity.
\end{proof}

\section{A direct graphon Moon--Moser inequality}

We next prove the clique-density inequalities needed to compare
Theorem~\ref{thm:gluing} with the polynomial certificate.

\begin{lemma}[Weighted Cauchy--Schwarz]\label{lem:weighted-cs}
If \(A\ge0\) and \(A,A\eta,A\eta^2\) are integrable, then
\[
    \left(\int A\eta\right)^2
    \le
    \left(\int A\right)\left(\int A\eta^2\right).
\]
\end{lemma}

\begin{proof}
This is Cauchy--Schwarz applied to \(\sqrt A\) and \(\sqrt A\,\eta\).
All functions below are bounded finite products, so the integrability
hypotheses hold automatically.
\end{proof}

Fix an integer \(j\ge2\).  On \(\Omega^{j-1}\), put
\[
    A(z_1,\dots,z_{j-1})
    :=
    \prod_{1\le a<b\le j-1}W(z_a,z_b)
\]
and
\[
    \eta(z_1,\dots,z_{j-1})
    :=
    \int_\Omega\prod_{a=1}^{j-1}W(x,z_a)\,d\mu(x).
\]
Then
\[
    \int A=t_{j-1},
    \qquad
    \int A\eta=t_j.
\]
Let
\[
    J_j:=\int A\eta^2.
\]
This is the homomorphism density of \(K_{j+1}\) with the edge between two
distinguished vertices deleted.  Lemma~\ref{lem:weighted-cs} gives
\begin{equation}\label{eq:clique-cs}
    t_j^2\le t_{j-1}J_j.
\end{equation}

\begin{lemma}[Finite cube inequality]\label{lem:cube}
Assign a number \(a_e\in[0,1]\) to every edge of the complete graph on
\(j+1\) labelled vertices.  Define
\[
    K:=\prod_ea_e,
    \qquad
    B:=\sum_e\prod_{f\ne e}a_f,
\]
and
\[
    S:=\sum_v
       \prod_{e\subseteq [j+1]\setminus\{v\}}a_e.
\]
Then
\begin{equation}\label{eq:cube}
    2B\le S+(j^2-1)K.
\end{equation}
\end{lemma}

\begin{proof}
The difference between the right- and left-hand sides is affine in each
variable \(a_e\) separately.  By induction over the finite edge set, a
multiaffine function that is nonnegative on all Boolean assignments is
nonnegative on the whole cube \([0,1]^E\): for one selected coordinate,
its value is the convex combination
\[
    (1-a_e)F(a_e=0)+a_eF(a_e=1).
\]

It remains to check Boolean assignments.  If at least two edge variables are
zero, then \(B=K=0\) and \(S\ge0\).  If exactly one edge variable is zero,
then \(B=1\), \(K=0\), and precisely the two deletions of its endpoints
contribute to \(S\), so \(S=2\).  If no variable is zero, then
\[
    B=\binom{j+1}{2},
    \qquad S=j+1,
    \qquad K=1,
\]
and both sides of \eqref{eq:cube} equal \(j(j+1)\).
\end{proof}

\begin{lemma}[Second-moment symmetrization]\label{lem:second-moment}
For every integer \(j\ge2\),
\begin{equation}\label{eq:second-moment}
    jJ_j\le t_j+(j-1)t_{j+1}.
\end{equation}
\end{lemma}

\begin{proof}
In Lemma~\ref{lem:cube}, substitute
\(a_{\{u,v\}}=W(x_u,x_v)\) and integrate over
\(\Omega^{j+1}\).  Every edge-deleted term has density \(J_j\), by a
permutation of the \(j+1\) coordinates; there are
\(\binom{j+1}{2}\) such terms.  Every vertex-deleted term has density \(t_j\),
and there are \(j+1\) of them.  The complete term has density \(t_{j+1}\).
Thus \eqref{eq:cube} integrates to
\[
    2\binom{j+1}{2}J_j
    \le
    (j+1)t_j+(j^2-1)t_{j+1}.
\]
Both sides have the positive factor \(j+1\); cancellation gives
\eqref{eq:second-moment}.
\end{proof}

\begin{theorem}[Graphon Moon--Moser recurrence]\label{thm:moon-moser}
For every graphon \(W\) and every integer \(j\ge2\),
\begin{equation}\label{eq:moon-moser}
    j t_j^2
    \le
    t_{j-1}t_j+(j-1)t_{j-1}t_{j+1}.
\end{equation}
\end{theorem}

\begin{proof}
Multiply \eqref{eq:clique-cs} by \(j\), then use
Lemma~\ref{lem:second-moment} and \(t_{j-1}\ge0\):
\[
\begin{aligned}
    jt_j^2
    &\le jt_{j-1}J_j\\
    &\le t_{j-1}\bigl(t_j+(j-1)t_{j+1}\bigr),
\end{aligned}
\]
which expands to \eqref{eq:moon-moser}.
\end{proof}

\section{Clique polynomials and denominator-free ratios}

For \(j\ge1\), define
\[
    c_j(p):=1-(j-1)(1-p).
\]
Then \(c_1=1\), \(c_2=p\), and
\[
    A_j(p)=\prod_{h=1}^{j}c_h(p).
\]

\begin{lemma}[Signs at the required threshold]\label{lem:factor-signs}
Assume \(r\ge3\) and \(p\ge1-1/(r-1)\).  Then
\[
    c_j(p)>0\quad(1\le j<r),
    \qquad
    c_r(p)\ge0.
\]
Consequently,
\[
    A_s(p)>0\quad(0\le s<r),
    \qquad
    A_r(p)\ge0.
\]
\end{lemma}

\begin{proof}
The density hypothesis is equivalent to
\(1-p\le1/(r-1)\).  If \(j<r\), then
\[
    (j-1)(1-p)
    \le
    \frac{j-1}{r-1}<1.
\]
For \(j=r\), the strict inequality becomes a weak one.  The product
statements follow.
\end{proof}

\begin{theorem}[Adjacent clique step]\label{thm:clique-step}
Assume \(r\ge3\) and \(p=t_2\ge1-1/(r-1)\).  For every
\(2\le j\le r\),
\begin{equation}\label{eq:clique-step}
    t_j\ge c_j(p)t_{j-1}.
\end{equation}
For every \(j<r\), one also has \(t_j>0\).
\end{theorem}

\begin{proof}
The base case is
\[
    t_2=p=c_2(p)t_1.
\]
It is positive because \(p\ge1-1/(r-1)>0\).

Suppose \(2\le j<r\), the estimate
\(t_j\ge c_jt_{j-1}\) is known, and \(t_j,t_{j-1}>0\).  Rearranging
Theorem~\ref{thm:moon-moser} gives
\begin{equation}\label{eq:mm-rearranged}
    (j-1)t_{j-1}t_{j+1}
    \ge
    t_j(jt_j-t_{j-1}).
\end{equation}
The induction hypothesis implies
\[
\begin{aligned}
    jt_j-t_{j-1}
    &\ge (jc_j-1)t_{j-1}\\
    &=(j-1)c_{j+1}t_{j-1},
\end{aligned}
\]
where the final equality is the elementary identity
\(jc_j-1=(j-1)c_{j+1}\).  Substitute this in
\eqref{eq:mm-rearranged}.  The factor \((j-1)t_{j-1}\) is strictly positive,
so cancellation yields
\[
    t_{j+1}\ge c_{j+1}t_j.
\]
If \(j+1<r\), Lemma~\ref{lem:factor-signs} makes the right-hand side
strictly positive.  This proves both claims by induction.  Notice that no
positivity of \(t_r\) is asserted at the boundary; the last coefficient can
be zero.
\end{proof}

\begin{corollary}[Clique polynomial and cross-ratio bounds]
Under the hypotheses of Theorem~\ref{thm:clique-step}, for \(0\le j\le r\),
\begin{equation}\label{eq:clique-poly-lower}
    t_j\ge A_j(p).
\end{equation}
Moreover, whenever \(0\le s\le j\le r\),
\begin{equation}\label{eq:cross-ratio}
    t_sA_j(p)\le t_jA_s(p).
\end{equation}
\end{corollary}

\begin{proof}
Iterate \eqref{eq:clique-step}, using \(t_0=t_1=A_0=A_1=1\), to obtain
\eqref{eq:clique-poly-lower}.  More generally, iteration from \(s\) to \(j\)
gives
\[
    t_j
    \ge
    t_s\prod_{h=s+1}^{j}c_h(p).
\]
Multiplying by \(A_s(p)\ge0\) and using
\(A_j=A_s\prod_{h=s+1}^{j}c_h\) proves
\eqref{eq:cross-ratio}.  For \(s=0\), the identical values at indices \(0\)
and \(1\) reduce the statement to the same calculation.  This proof never
divides by a possibly vanishing clique density or polynomial factor.
\end{proof}

\section{The certificate inequality}

\begin{theorem}[Pure clique-tree certificate bound]\label{thm:certificate}
Let \(H\) have a rooted \(r\)-pure clique-tree presentation, where \(r\ge3\).
For every graphon \(W\) with
\[
    p=t(K_2,W)\ge1-\frac1{r-1},
\]
one has
\begin{equation}\label{eq:certificate-bound}
    t(H,W)\ge B_H(p)
    =\prod_{i=0}^{m-1}\Tail_{s_i,r}(p).
\end{equation}
\end{theorem}

\begin{proof}
Set
\[
    S:=\prod_{i=0}^{m-1}t_{s_i},
    \qquad
    A:=\prod_{i=0}^{m-1}A_{s_i}(p).
\]
By Lemmas~\ref{lem:separator-size} and \ref{lem:factor-signs}, every
\(A_{s_i}(p)>0\).  By \eqref{eq:clique-poly-lower},
\(t_{s_i}\ge A_{s_i}(p)>0\).  Hence
\begin{equation}\label{eq:SA-positive}
    S>0,
    \qquad A>0.
\end{equation}

Apply \eqref{eq:cross-ratio} with \(s=s_i\) and \(j=r\), then multiply the
\(m\) resulting nonnegative inequalities.  This gives
\begin{equation}\label{eq:ratio-product}
    S A_r(p)^m
    \le
    t_r^m A.
\end{equation}
The gluing theorem gives \(t_r^m\le t(H,W)S\).  Since \(A\ge0\),
\[
    S A_r(p)^m
    \le
    t_r^mA
    \le
    t(H,W)SA.
\]
Cancel the strictly positive factor \(S\) to obtain
\begin{equation}\label{eq:cross-certificate}
    A_r(p)^m\le t(H,W)A.
\end{equation}
By \eqref{eq:certificate-cross}, the left-hand side is \(AB_H(p)\).
Cancel the strictly positive factor \(A\) from
\eqref{eq:cross-certificate}.  The result is
\(B_H(p)\le t(H,W)\).

This argument includes the exact threshold
\(p=1-1/(r-1)\), where \(A_r(p)=0\), and \(p=1\).  No division by \(A_r\)
or \(t_r\) occurs.
\end{proof}

\begin{proof}[Proof of Theorem~\ref{thm:main}]
Choose a rooted maximal-clique-tree presentation.  Purity makes every bag
have size \(r\), so Theorem~\ref{thm:certificate} applies.  By
Lemma~\ref{lem:chromatic}, \(\chi(H)=r\), and hence its density hypothesis is
exactly the required interval
\[
    p\ge1-\frac1{\chi(H)-1}.
\]
For \(p<1\), Lemma~\ref{lem:target} identifies the certificate with the
stated chromatic expression.  At \(p=1\), the certificate equals \(1\), so
Theorem~\ref{thm:certificate} gives \(t(H,W)\ge1\); the universal bound
\(t(H,W)\le1\) gives equality.
\end{proof}

\section{Concrete consequences and checks}

\begin{corollary}[Generalized books]
Let \(B_{r,m}=K_{r-1}\vee\overline{K_m}\), namely \(m\) copies of \(K_r\)
sharing one \(K_{r-1}\).  Then, for
\(p\ge1-1/(r-1)\),
\[
    t(B_{r,m},W)
    \ge
    A_{r-1}(p)\bigl(1-(r-1)(1-p)\bigr)^m.
\]
For \(r=3\), this is
\[
    t(B_{3,m},W)\ge p(2p-1)^m,
    \qquad p\ge\frac12.
\]
\end{corollary}

\begin{proof}
Use one root bag and \(m-1\) non-root bags, each with separator size \(r-1\).
The certificate is
\[
    A_r(p)
    \left(\frac{A_r(p)}{A_{r-1}(p)}\right)^{m-1}
    =
    A_{r-1}(p)c_r(p)^m.
\]
The displayed quotient is only shorthand for the explicit tail factor
\(c_r(p)\); the proof itself uses no division.
\end{proof}

\begin{corollary}[Clique windmills]
Let \(F_{r,s,m}\) consist of \(m\) copies of \(K_r\), all sharing the same
\(K_s\), with no other overlap.  Then
\[
    t(F_{r,s,m},W)
    \ge
    A_s(p)\left(\frac{A_r(p)}{A_s(p)}\right)^m
\]
for \(p\ge1-1/(r-1)\), where every quotient again denotes the polynomial
tail \(\Tail_{s,r}(p)\).  For \(r=3,s=1\), this is the friendship bound
\[
    t(F_{3,1,m},W)
    \ge
    \bigl[p(2p-1)\bigr]^m,
    \qquad p\ge\frac12.
\]
\end{corollary}

\begin{corollary}[Pure triangle trees]
Let \(H\) be connected and suppose all maximal cliques are triangles.  If
there are \(m\) triangle bags and exactly \(a\) non-root separators are
edges, then
\[
    t(H,W)
    \ge
    p^{m-a}(2p-1)^m,
    \qquad p\ge\frac12.
\]
\end{corollary}

\begin{proof}
Here
\[
    A_3(p)=p(2p-1),
    \qquad
    \Tail_{2,3}(p)=2p-1,
    \qquad
    \Tail_{1,3}(p)=p(2p-1).
\]
The root contributes \(A_3\), the \(a\) edge separators contribute
\(2p-1\), and the remaining \(m-1-a\) separators have size \(1\) and
contribute \(p(2p-1)\).  Multiplication gives the formula.
\end{proof}

The one-bag case reduces to \(t_r\ge A_r(p)\).  With two bags glued along
\(K_s\), Theorem~\ref{thm:gluing} reduces to
\[
    t(K_r\cup_{K_s}K_r,W)t_s\ge t_r^2,
\]
which can independently be checked by weighted Cauchy--Schwarz.  These are
useful regression cases for both the paper proof and its formalization.

\section{Correspondence with the Lean verification}

The directory \path{lean/Taeyoung/Methods/PureChordal} contains a Lean
4/Mathlib proof with
no \texttt{sorry}, \texttt{admit}, declared custom axiom, or
\texttt{native\_decide}.  The mathematical correspondence is as follows.

The important standard supplied by that project is the public theorem shape.
In the eventual integrated project described in
\texttt{LEAN\_VERIFICATION\_ARCHITECTURE.md}, the shared foundation should
define \texttt{chromaticTarget} with its polynomial value at \(p=1\), and the
method-level file \texttt{Methods/PureChordal/Main.lean} should export a
statement of the following form.

\begin{verbatim}
theorem pureChordal_satisfiesLowerBound
    (hchordal : IsChordal H)
    (hpure : HasPureMaximalCliques H r)
    (hr : 3 <= r) :
    SatisfiesLowerBound H
\end{verbatim}

After unfolding the shared predicate, this means: for every probability space
and every graphon on it, the admissible edge-density hypothesis implies
\texttt{chromaticTarget H p <= homDensity H W}.  Internal certificate,
junction-density, regularization, and Moon--Moser lemmas need not be visible to
an Atlas example.  An example file should import only this method's
\texttt{Main.lean}, prove the graph is pure chordal, and apply the theorem.

\begin{center}
\begin{tabular}{p{0.39\textwidth}p{0.53\textwidth}}
\hline
This paper & Lean declaration or module \\
\hline
Definition~\ref{def:certificate} &
\texttt{PureCliqueTreeDecomp} in \texttt{Certificate.lean} \\
Lemmas~\ref{lem:separator-size}--\ref{lem:edge-multiplicity} &
\texttt{Certificate.lean} and \texttt{CliqueTreeCombinatorics.lean} \\
Lemma~\ref{lem:chromatic} &
\texttt{properAssignmentCount\_eq\_product} and
\texttt{chromaticPolynomial} in \texttt{ChromaticFactorization.lean} \\
Lemma~\ref{lem:target} &
\texttt{certificateBound\_eq\_eval\_chromaticPolynomial} \\
Lemma~\ref{lem:gibbs} &
\texttt{integral\_mul\_log\_div\_nonpos} in \texttt{Gibbs.lean} \\
Lemmas~\ref{lem:equal-marginal}--\ref{lem:constant-ratio} &
\texttt{CliqueMarginals.lean}, \texttt{JunctionDensity.lean}, and
\texttt{EntropyGluing.lean} \\
Theorem~\ref{thm:gluing} &
\texttt{homDensity\_mul\_sep\_ge\_cliqueDensity\_pow} \\
Regularization estimate \eqref{eq:regularization-error} &
\texttt{abs\_homDensity\_regularize\_sub\_le} in
\texttt{Regularization.lean} \\
Theorem~\ref{thm:moon-moser} &
\texttt{cliqueDensity\_moonMoser\_succ} in \texttt{MoonMoser.lean} \\
Theorem~\ref{thm:clique-step} and \eqref{eq:cross-ratio} &
\texttt{cliqueDensity\_step\_lower} and
\texttt{cliqueDensity\_mul\_cliquePoly\_le} \\
Theorem~\ref{thm:certificate} &
\texttt{certificateBound\_le\_homDensity} \\
Theorem~\ref{thm:main} &
\texttt{pureChordal\_chromaticPolynomial\_lower\_bound} in
\texttt{Main.lean} \\
Common catalogue proposition &
\texttt{pureChordal\_satisfiesLowerBound} in \texttt{Main.lean} \\
\hline
\end{tabular}
\end{center}

The low-level Lean chromatic-expression theorem assumes \(p\ne1\), exactly
because the literal real expression \(1/(1-p)\) is undefined at \(p=1\).
The Lean certificate theorem has no such restriction and proves the
endpoint value directly.  The catalogue-facing theorem
\texttt{pureChordal\_satisfiesLowerBound} performs this split explicitly:
it uses the certificate/polynomial identity only for \(p\ne1\), and proves
the value (1) at (p=1) from the certificate bound.  Thus no division by
zero is hidden in the common formal interface.

The soundness module
\path{lean/Taeyoung/Methods/PureChordal/CheckComplete.lean} prints the axiom
dependencies of the two user-facing theorems.  The reported list is
\[
    \{\texttt{propext},\ \texttt{Classical.choice},\
      \texttt{Quot.sound}\},
\]
the standard Mathlib classical axioms, and contains no \texttt{sorryAx}.

\section{Dependency and hidden-assumption audit}

The proof uses the following ingredients, all proved above or isolated in the
Lean development:
\begin{enumerate}
    \item the maximal-clique-tree characterization of finite chordal graphs;
    \item finite set, finite product, and finite colouring identities;
    \item finite-product measure invariance under coordinate relabeling;
    \item Tonelli/Fubini for bounded nonnegative functions;
    \item \(\log u\le u-1\) for \(u>0\);
    \item Cauchy--Schwarz and the elementary multiaffine cube lemma;
    \item the explicit regularization error
    \(|t(F,W_\varepsilon)-t(F,W)|\le e(F)\varepsilon\).
\end{enumerate}

It does not assume that \(W\) is regular, a finite step graphon, or a limit of
finite graphs.  It never replaces homomorphism density by injective density.
It uses no lower bound on the pointwise degree function and no minimum-degree
hypothesis.  Purity is essential in Lemma~\ref{lem:equal-marginal}: uniform
clique laws of different bag sizes need not have compatible separator
marginals.  Therefore this proof does not establish the theorem for chordal
graphs with mixed maximal-clique sizes.

\begin{thebibliography}{9}

\bibitem{MoonMoser}
J. W. Moon and L. Moser,
\emph{On a problem of Tur\'{a}n},
Magyar Tud. Akad. Mat. Kutat\'{o} Int. K\"ozl. \textbf{7} (1962), 283--286.

\bibitem{Agnarsson}
G. Agnarsson,
\emph{On chordal graphs and their chromatic polynomials},
Mathematica Scandinavica \textbf{93} (2003), 240--246.

\end{thebibliography}

\end{document}
